% Authored. Numbers come from generated/macros.tex.

\section{Introduction}
\label{sec:introduction}

A country's fiscal position is normally reported as one number. Portugal recorded a
general-government balance of \LatestGg\ million euro in \LatestYear, or
\LatestGgPct\% of GDP, and it is that figure which enters fiscal rules, credit
assessments and public debate.

The number is an aggregate of institutionally distinct parts. Under ESA 2010 the
general government sector S.13 is the consolidation of central government
(S.1311), state government (S.1312), local government (S.1313) and social security
funds (S.1314) \citep{eurostat_esa2010}. These are not administrative
subdivisions of one budget. They have different revenue bases --- general taxation
in one case, earmarked social contributions in another --- different expenditure
mandates, and different legal constraints on borrowing. Their balances are
therefore determined by different things, and there is no reason to expect them to
move together.

Aggregation hides that. In \LatestYear\ Portugal's positive aggregate balance was
the sum of a Central Government balance of \LatestCentral\ million euro, a Regional
and Local balance of \LatestRegional\ million, and a Social Security Funds balance
of \LatestSsf\ million. The headline figure and the composition behind it support
quite different descriptions of the same year, and only the headline figure is
routinely reported.

\subsection{Research questions}

This paper addresses two questions.

\begin{enumerate}
\item \textbf{How has Portugal's general-government balance been composed across
institutional subsectors since \PanelStart, and how persistent are the
contributions of Central Government, Regional and Local Government, and Social
Security Funds?}
\item \textbf{What accounting mechanisms account for the changing contribution of
the Social Security Funds, particularly in the recent period?}
\end{enumerate}

Both are descriptive questions about recorded quantities, and we answer them as
such. The central object is an accounting identity,
\[
\Bgg_t = \Bc_t + \Brl_t + \Bssf_t,
\]
which holds by construction. Decomposing it reallocates an observed number; it
estimates nothing and identifies nothing. We are explicit about this throughout,
because the decomposition invites a counterfactual reading --- that the aggregate
balance would have been different under a different institutional arrangement ---
which the identity cannot support.

\subsection{Related work and where this paper sits}

Three strands of existing work bear on these questions, and each leaves a
different part of them open.

The first is the statistical framework itself. ESA 2010 defines the subsectors and
requires member states to report deficit and debt for each of them, and Eurostat
and the ECB publish the resulting series
\citep{eurostat_esa2010,eurostat_ssf_glossary,eurostat_gov_main,ecb_gfs}. The
composition of the aggregate balance is therefore not hidden: it is published. What
the framework provides is a definition and a data supply, not an analysis of how a
particular country's composition has behaved over time.

The second is the literature on Portuguese fiscal sustainability, which works
predominantly at the aggregate level and over long horizons.
\citet{marinheiro2006} tests the sustainability of Portuguese fiscal policy on more
than a century of annual data, using stationarity and cointegration tests on
aggregate revenue, expenditure and debt. That question --- whether the aggregate
path satisfies a solvency condition --- is a different question from ours, and it is
answered without decomposing the aggregate.

The third concerns Portuguese Social Security financing specifically.
\citet{cardoso2019} examines the equity and sustainability of the system's
financing, noting that social-protection entitlements follow a dynamic that at
times diverges from the cyclical path of the revenue that funds them; the
Portuguese Public Finance Council has also examined the Social Security Financial
Stabilisation Fund in the same context \citep{cfp_fefss}. This literature is
concerned with the system's own finances and its long-run adequacy, rather than
with the arithmetic by which that system's balance enters the national aggregate.

A fourth strand is relevant less as a precedent than as a caution. Work on
stock-flow adjustments has shown that the reported balance is an incomplete account
of what happens to debt: \citet{vonhagen2006} find that debt accumulates by more
than accumulated deficits and read the residual as evidence of accounting responses
to fiscal rules, while \citet{seiferling2013} argues on integrated public-finance
data that the residuals are smaller than previously assumed and are not well
explained by fiscal transparency. We take no position in that debate. We take from
it the discipline that a balance is a recorded accounting quantity whose relation
to anything else has to be demonstrated rather than assumed --- which is why this
paper stops at the accounting and says so.

To our knowledge no study assembles and analyses the subsector composition of the
Portuguese general-government balance systematically over a window of this length.
We state that as a gap in the work we have found rather than as a claim of
priority: the underlying series are published, and any competent user of the
Eurostat or CFP data could construct parts of what follows. The contribution is in
the assembly, the harmonisation, the discipline applied to the resulting statistics
and the reproducibility of the whole, not in access to anything unavailable.

\subsection{Contribution}

We make six contributions.

\begin{enumerate}
\item \textbf{A harmonised long-run panel.} An annual panel of net lending and net
borrowing for General Government and its three subsectors covering
\PanelStart--\PanelEnd, built by splicing the Banco de Portugal / INE historical
long series to the modern national accounts. The \textbf{1995} methodological
boundary is retained as a documented splice, both vintages of the overlapping year
are kept, and nothing is chained, smoothed or calibrated across it.

\item \textbf{An exact decomposition, in levels and in changes.} The aggregate
balance and each of its annual changes are decomposed into subsector components
that close to the sources' rounding. Changes are additionally reported scaled by
current-year GDP, which makes episodes forty years apart comparable without
turning the decomposition into a ratio that no longer adds up.

\item \textbf{Persistence stated per statistical regime.} Sign frequencies, run
lengths and one-year sign transitions, with magnitudes computed inside each regime
rather than pooled across the splice. Pooling is not a presentational choice here:
the aggregate balance averages \GgMeanHistorical\% of GDP before 1995 and
\GgMeanModern\% after, and the pooled \GgMeanPooled\% describes neither period.

\item \textbf{Attribution of the largest episodes to revenue and expenditure.} The
largest annual improvements and deteriorations, each attributed to a subsector and
then to the revenue and expenditure movements that composed it.

\item \textbf{Two Social Security boundaries, separated and quantified.} The
ESA 2010 Social Security Funds sector and the budget-execution systems published by
the CFP are different accounting objects. We report them side by side, never add
them, and measure the gap between them --- which is non-zero in every one of the
\SsfBoundaryYears\ years they overlap.

\item \textbf{A validation layer that separates two different tests.} Identity
closure asks whether an extraction is arithmetically self-consistent; source
agreement asks whether two independently published sources report the same number.
The first cannot establish the second. Here the identities close to
\MaxClosureResidual\ million euro while two published sources still disagree by
\MaxSourceDisagreement\ million euro on one subsector-year.
\end{enumerate}

\subsection{Principal findings}

The composition is asymmetric and it is stable. Over \PanelYears\ years the
aggregate balance is positive in \NumPositiveYears\
(\PositiveYearList). The Central Government balance is negative in all
\CentralNegativeYears\ observations. The Social Security balance is positive in
\SsfPositiveYears, with a longest unbroken run of \SsfLongestPositiveRun\ years.
Regional and Local Government is the only subsector that changes sign with any
regularity, and it is small in both directions.

In each of the \NumPositiveYears\ years with a positive aggregate balance, the
combined non-Social-Security balance is negative and the Social Security balance
exceeds it in absolute size. In \LatestYear\ the non-Social-Security balance was
\LatestNonSsf\ million euro against a Social Security balance of \LatestSsf\
million, an offset ratio of \LatestOffset. This is an arithmetic relation between
two recorded numbers and we do not present it as anything else.

Magnitudes are regime-dependent and must be read that way. The Central Government
balance averages \CentralMeanHistorical\% of GDP over \PanelStart--1994 and
\CentralMeanModern\% over 1995--\PanelEnd; the Social Security balance rises from
\SsfMeanHistorical\% to \SsfMeanModern\%.

The headline sign is not the underlying sign. Central Government records a negative
balance in every year of the panel, but its primary balance --- the balance
excluding interest --- is positive in \PrimaryPositiveYears\ of the
\PrimaryObservedYears\ years for which the detailed accounts exist, from
\PrimaryPositiveFirst\ to \PrimaryPositiveLast. Interest averaged
\InterestMeanPct\% of GDP over that panel and peaked at \InterestPeakPct\%. Both
statements are descriptive and they are not in conflict; interest is the arithmetic
that separates them. A conclusion that Central Government runs a permanent
underlying deficit does not follow from the headline series.

Finally, the recent movement in the Social Security position is concentrated in one
of the two internal budget systems. The Previdential balance moves from
\PrevidentialFirst\ million euro in \PrevidentialFirstYear\ to
\PrevidentialLatest\ million in \LatestYear, while the Citizenship balance stays
small in both directions and stands at \CitizenshipLatest\ million. That is a
description of the published series and not an account of what caused it.

\subsection{Roadmap}

Section~\ref{sec:framework} defines the institutional and accounting objects, and
in particular why Central Government is not a synonym for the State in the
budgetary sense. Section~\ref{sec:data} describes the sources, the harmonisation,
the 1995 splice, the vintages and the cross-source validation.
Section~\ref{sec:method} states the decompositions and the persistence measures.
Section~\ref{sec:results} presents the results. Section~\ref{sec:robustness} sets
out the accounting boundaries and the robustness of the weaker statistics.
Section~\ref{sec:conclusion} concludes with what the evidence does and does not
establish. Appendix~\ref{sec:changepoints} reports the change-point sensitivity in
full and Appendix~\ref{sec:reproducibility} the reproduction path.
